\documentclass[10pt,a4paper]{amsart}
\usepackage[margin=19mm]{geometry}
\usepackage{amsmath,amssymb,mathtools}
\usepackage[scr=rsfs,cal=euler]{mathalfa}
\usepackage{xfrac}
\usepackage[unicode,hidelinks,bookmarks=false]{hyperref}
\newcommand{\ie}{i.e.\ }
\newcommand{\cf}{cf.\ }
\newcommand{\resp}{respectively\ }
\newcommand{\Subsection}{\subsection}
\providecommand{\nobreakdash}{\nobreak\hbox{-}}
\setlength{\emergencystretch}{2em}
\multlinegap0pt
\usepackage{xcolor}
\usepackage{amssymb}
\usepackage{mathtools}
\usepackage{enumerate}
\newtheorem{proposition}{Proposition}
\newcommand{\A}{\mathcal{A}}
\newcommand{\B}{\mathcal{B}}
\newcommand{\Z}{\mathbb{Z}}
\newcommand{\hilb}{\mathcal{H}}
\title{\huge{Discrete $C^*$-Frobenius Algebras and Infinite-Index Extensions in Algebraic Quantum Field Theories}}
\author{{\sc ZIYUN XU}}
\date{Source: arXiv:2609.14961v1 (CC BY 4.0)}
\begin{document}
\maketitle

\noindent\emph{Source and licence.} \huge{Discrete $C^*$-Frobenius Algebras and Infinite-Index Extensions in Algebraic Quantum Field Theories}. Version arXiv:2609.14961v1 (CC BY 4.0), distributed by arXiv under the Creative Commons Attribution 4.0 International licence, http://creativecommons.org/licenses/by/4.0/, as recorded in the arXiv licence metadata for this exact version and shown on the abstract page https://arxiv.org/abs/2609.14961v1. Full licence notice: \texttt{licenses/arxiv-2609.14961-CC-BY-4.0.txt}.\par\smallskip

\noindent\textit{Editorial scope.} Counted unit: the article's proposition extension (source0.tex lines 911--914). The article's complete original proof of it is delivered with it (source0.tex lines 915--931, one \texttt{proof} environment of 17 lines). No mathematics is written, repaired or re-derived: the statement and the proof are the article's own lines, in the article's own notation, and no retained line is substituted or re-flowed.  No further source-local context is needed: every cross-reference of every retained line resolves inside the two delivered spans.  Nothing was omitted from any retained span, so the package carries no omission record.  No retained line contains a citation, so the wrapper carries no bibliography and no bibliography entry.  No retained line contains a figure environment, an image-inclusion command or drawing code, so no image is delivered and the package declares no asset dependency.  Editorial formatting only, in addition: an AMS \texttt{amsart} preamble that restates the article's own declarations and macros used by the retained lines, namely the environments \texttt{proposition} on separate counters, together with the standard packages those lines need.  The retained environments print in this document as Proposition 1 in order of appearance, whereas the article prints the same items with the numbering of its own full document.  The complete native source of the article as distributed in the arXiv e-print for this exact version is bundled unchanged as \texttt{sources/210379/source0.tex}.
\begin{proposition}\label{extension}
For any $\widetilde{I}, \widetilde{J} \in\widetilde{\mathcal{I}}$, we have 
    $ \pi_{\hilb_\bullet, I \times J}\left(\A(I)\otimes\A^\mathfrak{r}(J)\right) \subset  \B^{\mathfrak{L}}(\widetilde{I}\times\widetilde{J}) \cap \B^{\mathfrak{R}}(\widetilde{I}\times\widetilde{J})  $. 
\end{proposition}

\begin{proof}
Fix a vector $\xi\in\hilb_\bullet$ such that $\operatorname{supp}\xi = \left\{ i\in\Z: \widetilde{T_{i}}\xi \neq 0 \right\}$ is a finite subset of $\Z$ and 
$\xi =\sum_{j\in \operatorname{supp}\xi} \sum_{k =1}^{\xi(j)}\eta_j^{(j_k)} \otimes \zeta_j^{(j_k)}$,
where for any $j\in\operatorname{supp}\xi$,
$\xi(j)\in\Z_{>0}$ and for any $1 \leq k \leq \xi(j)$,
$\eta_j^{(j_k)}, \zeta_j^{(j_k)} \in \hilb_j$.
For $\widetilde{I}, \widetilde{J} \in\widetilde{\mathcal{I}}$, $x\in \A(I)$ and $y \in \A^\mathfrak{r}(J)$, we have: 
    \begin{align*}
        \mu_0^{\mathfrak{L}}(L(x\Omega, \widetilde{I})\otimes \Theta L(y\Omega, \mathfrak{r}\widetilde{J})\Theta^{-1})\xi&=  \sum_{j\in\operatorname{supp}\xi}\sum_{k =1}^{\xi(j)}\widetilde{T_{j}}^*\widetilde{V^{j}_{0,j}}[(L(x\Omega, \widetilde{I})\eta_j^{(j_k)})\otimes (\Theta L(y\Omega, \mathfrak{r}\widetilde{J})\Theta^{-1}\Theta_j\zeta_j^{(j_k)})] \\
        & = \sum_{j\in\operatorname{supp}\xi}\sum_{k =1}^{\xi(j)}\widetilde{T_{j}}^*(V^j_{0,j}L(x\Omega, \widetilde{I})\eta_j^{(j_k)})\otimes(\overline{V^j_{0,j}}\Theta L(y\Omega, \mathfrak{r}\widetilde{J})\zeta_j^{(j_k)})\\
        & = \sum_{j\in\operatorname{supp}\xi}\sum_{k =1}^{\xi(j)}\widetilde{T_{j}}^* (V^j_{0,j}L(x\Omega, \widetilde{I})\eta_j^{(j_k)})\otimes (\Theta_jV^j_{0,j}(L(y\Omega, \mathfrak{r}\widetilde{J})\zeta_j^{(j_k)})) \\
        & = \sum_{j\in\operatorname{supp}\xi}\sum_{k =1}^{\xi(j)}\widetilde{T_{j}}^*  (\pi_{j,I}(x)\eta_j^{(j_k)} \otimes \Theta_j \cdot \pi_{j,\mathfrak{r}J}(y) \cdot \Theta_j^{-1} \cdot \Theta_j \zeta_j^{(j_k)}) \\
        & =  \left(\sum_{j\in\operatorname{supp}\xi}\widetilde{T_{j}}^* (\pi_{j,I}(x) \otimes \pi_{j,J}^{\mathfrak{r}}(y))\widetilde{T_{j}}\right)\xi \\
        &= \pi_{\hilb_\bullet, I\times J}(x\otimes y)\xi.
    \end{align*}
    By a  similar argument, we obtain $\pi_{\hilb_\bullet, I \times J}(\A(I)\otimes\A^\mathfrak{r}(J)) \subset  \B^{\mathfrak{R}}(\widetilde{I}\times\widetilde{J}) $ .
\end{proof}

\end{document}
